\section{What richer observations recover}
Let $X$ be the ordered physical state space. A disclosed experimental context
$y$ draws a state with probabilities $K(y,x)$, so $K$ is row stochastic.
Commitment observations of a table $f$ are $Kf$. These are conditional-state
lotteries, not a forward signal channel without a prior. The nullspace of
$K$ describes exactly the mean information lost. A functional of $f$ is
identified precisely when its coefficient vector lies in the row span of $K$.
These are standard linear inverse-problem facts.

\begin{proposition}[Mean versus consequence-law recovery; P8-O1/O5]
The design identifies every real payoff table iff $\ker K=0$. On a common
finite consequence set with at least two ordered values, observations under
every common strictly increasing expected-utility representation identify
every deterministic consequence table iff
\[
 \ker K\cap\{-1,0,1\}^{X}=\{0\}.
\]
\end{proposition}
\begin{proof}
Subtract mean records for the first assertion. Robust equality of consequence
lotteries is equality of each atom probability, so for tables $f,h$ it says
$K(\mathbf1_{f=c}-\mathbf1_{h=c})=0$ for each consequence $c$.
Every indicator difference is ternary. Conversely a nonzero ternary kernel
vector is the difference of the indicators of its positive and negative
supports, giving two indistinguishable binary consequence tables.
\end{proof}
For $K=(1/3,2/3)$ means do not identify arbitrary tables, but its kernel
contains no nonzero ternary vector. Observing the whole consequence lottery
identifies both deterministic entries by their different weights. Conditional
stochastic consequences do not have this ternary restriction.

The uniform pair threshold is also sensitive to known bias. For a
transposition $T$ and a known probability $q\in[0,1]\setminus\{1/2\}$,
\[
 (qI+(1-q)T)^{-1}=\frac{qI-(1-q)T}{2q-1}.
\]
One biased pair mask, at every reference state, therefore identifies all mean
entries. On a nonzero antisymmetric space the inverse amplifies worst-case
Euclidean noise by $1/|2q-1|$. This is exact identification, with signed
reconstruction coefficients, rather than practical perfect disclosure.

\paragraph{The missing simultaneous information.}
Suppose a laboratory logs one controller protocol
$g(s)_{p(i)}=f_i(s_i)$, with each local button correspondence calibrated once
across all recipients and trials. Observe each recipient's conditional receipt
under this same joint intervention, rather than choosing different opponent
correspondences in separate focal comparisons. For discrepancy arrays
$d_i(s)=u_{p(i)}(gs)-u_i(s)$, the records are $Kd_i$.
If $K$ has trivial kernel, vanishing discrepancies are equivalent to $g$
being an exact automorphism. For a specified game it suffices that no nonzero
attainable discrepancy tuple lies in the channel kernels. The laboratory
protocol is physical data; agreement of optimized assignment values alone
does not establish it. Joint receipt distributions are unnecessary for
deterministic games: all recipient means conditional on the same logged
intervention suffice. Player transitivity additionally requires identified
protocols connecting every role.

\paragraph{Minimal and noisy symmetry tests.}
Suppose the unknown game belongs to a disclosed linear utility class $W$,
and write $D_g=R_g-I$ for the discrepancy operator of the calibrated protocol.
Observations $Lu$ identify whether $g$ is a symmetry throughout $W$ iff
\[
 \ker(L|_W)\subseteq\ker(D_g|_W).
\]
Necessity follows by comparing an invisible perturbation with the zero game;
sufficiency follows because $D_g|_W$ then factors through $L|_W$.
Consequently at least $\operatorname{rank}(D_g|_W)$ independent scalar
linear observations are necessary, and unrestricted linear designs attain
this bound. Restricted physical contexts may require more. This can need
less information than full utility recovery. On the full class, observe
all but one discrepancy entry in each cycle of $R_g$; each omitted entry
is determined by its cycle's zero sum.

For discrepancy observations $y=Md+e$, $d\in\Delta=D_gW$, assume
$\|e\|_2\le\eta$ and $\Delta\ne0$. Put
$\alpha=\min_{d\in\Delta,\|d\|_2=1}\|Md\|_2$. If $\alpha>0$,
\[
 \|d\|_2\le (\|y\|_2+\eta)/\alpha.
\]
This follows from the triangle inequality and the restricted minimum
singular value. Under the earlier common calibration and matched transports,
a sup-norm payoff discrepancy at most $\epsilon$ gives each matched source/target
role-value gap at most $\epsilon$ for every common belief. Since
$\|d\|_\infty\le\|d\|_2$, the observation bound supplies such a certificate.
With positive noise one cannot uniformly distinguish exact zero from every
nonzero discrepancy in an unrestricted nonzero $\Delta$. Against the
separated alternative $\|d\|_2\ge\epsilon$ with $\epsilon>0$, a deterministic uniform test
exists precisely when $\alpha\epsilon>2\eta$: the null and alternative
error balls are disjoint exactly then. Approximate payoff preservation does
not imply an exact automorphism, and fixed-tolerance approximate maps need
not form a group. These are standard linear estimation bounds applied to
specified interventions, not uncalibrated interpersonal welfare comparisons.

\paragraph{Utility representations.}
Under that same protocol, preservation of all comparisons of profile lotteries
for every recipient characterizes cardinal automorphisms by finite
expected-utility uniqueness; preservation of all pure comparisons characterizes
ordinal automorphisms by finite weak-order equivalence. These observations are
within-recipient preference tests, not cross-role welfare equality. Alternatively,
full-rank mean tests on minmax-normalized tables identify cardinal maps, and on
rank tables identify ordinal maps. Normalize before observing: averaging does
not generally commute with nonlinear recoding. Such tests are invariant under
independent affine or ordinal recodings, respectively, but their numerical
calibration does not become the participant's money scale by convention.
